% !PRINTED 620
\weberpage{620}
%% [page furniture, not body text -- source-confirmed 2026-08-10]
%% 620 \quad Dreiundzwanzigster Abschnitt. \quad \S\,169.

%<<CONT
und aus (4) folgt:
\weq{7}{\alpha' \equiv \pm 1 \pmod{\mathfrak{m}};}
folglich ist $A'$ in $A$ enthalten, und nach \S\,164, 11. ist also $\mathfrak{K}(A)$ in $\mathfrak{K}(A')$ enthalten. Da nun nach \S\,158 jeder Teilungsk\"orper sich aus solchen zusammensetzen l\"a\ss t, deren Teiler ein Primideal oder eine Potenz eines Primideals ist, so folgt:

6. \gesperrt{Der Teilungsk\"orper der elliptischen Funktionen ist zur\"uckf\"uhrbar auf Ordnungsk\"orper und Kreisteilungsk\"orper}\footnote{Auf anderem Wege ist ein Teil dieses Satzes bewiesen in der Stra\ss burger Dissertation von \gesperrt{Daniel Bauer}, \glqq{}\"Uber den Teilungsk\"orper der elliptischen Funktionen mit singul\"arem Modul\grqq{}. Stra\ss burg 1903.}.
